% best_solution_lru_scheduler_improved.tex
\documentclass[11pt,a4paper]{article}

%% Typography & layout
\usepackage[T1]{fontenc}
\usepackage[utf8]{inputenc}
\usepackage{microtype}
\usepackage{geometry}
\geometry{margin=1in}
\setlength{\parskip}{0.9ex}
\setlength{\parindent}{0pt}

%% Packages for visuals and code
\usepackage{amsmath,amssymb}
\usepackage{booktabs}
\usepackage{enumitem}
\usepackage{titlesec}
\usepackage{fancyhdr}
\usepackage{tcolorbox}
\usepackage{hyperref}
\usepackage{caption}
\usepackage{listings}
\usepackage{xcolor}

%% Color palette (subtle)
\definecolor{accent}{HTML}{0B6FA4}
\definecolor{dark}{HTML}{222222}
\definecolor{shade}{HTML}{F6F9FC}

%% Header / footer
\pagestyle{fancy}
\fancyhf{}
\lhead{\textsc{Data Structures \& Systems Design — Assignment}}
\rhead{\textsc{Naman Gautam}}
\cfoot{\thepage}

%% Section title style
\titleformat{\section}{\large\bfseries}{\thesection}{1em}{}
\titleformat{\subsection}{\normalsize\bfseries}{\thesubsection}{1em}{}
\titleformat{\subsubsection}{\normalsize\itshape}{\thesubsubsection}{1em}{}

%% listings configuration
\lstset{
  basicstyle=\ttfamily\footnotesize,
  breaklines=true,
  numbers=left,
  numberstyle=\tiny,
  frame=single,
  rulecolor=\color{gray},
  backgroundcolor=\color{shade},
  xleftmargin=4pt,
  captionpos=b
}

%% Tcolorbox styles
\tcbset{
  mybox/.style={
    colback=shade, colframe=accent, boxrule=0.6pt, arc=3pt, leftrule=0.6mm, rightrule=0.6mm,
    top=2pt, bottom=2pt, left=4pt, right=4pt
  },
  summary/.style={mybox, fonttitle=\bfseries}
}

%% Document metadata
\title{\vspace{-12pt}\textbf{LRU Cache \& Event Scheduler} \\ \large Production-grade solutions — algorithms, analysis, and code}
\author{Naman Gautam \\ SE Intern Submission}
\date{\today}

\begin{document}
\maketitle

\begin{tcolorbox}[summary,title=Executive summary]
This report contains production-quality solutions for two tasks:
\begin{enumerate}[noitemsep,leftmargin=*]
  \item An \textbf{O(1)} LRU cache implemented using \texttt{std::unordered\_map} + \texttt{std::list}, with rationale and complexity analysis.
  \item An event scheduler with functions to test attendability, compute minimum rooms required, and assign concrete room indices using a heap-based greedy algorithm.
\end{enumerate}
Appendices include full, compilable C++ source files and a minimal test harness.
\end{tcolorbox}

\vspace{4pt}
\tableofcontents
\thispagestyle{empty}
\newpage

\section{Problem statements (concise)}
\subsection{LRU Cache}
Implement a fixed-capacity cache that supports:
\begin{itemize}
  \item \texttt{int get(int key)} — returns the value or \texttt{-1} if missing; marks the key most-recently-used.
  \item \texttt{void put(int key,int value)} — inserts or updates; evicts the least-recently-used item if capacity is exceeded.
\end{itemize}
Both operations must be amortized \(O(1)\).

\subsection{Event Scheduler}
Given a list of intervals \((\text{start},\text{end})\):
\begin{itemize}
  \item \texttt{can\_attend\_all(events)}: return \texttt{true} if no overlapping intervals exist (adjacent allowed).
  \item \texttt{min\_rooms\_required(events)}: compute the minimum number of rooms required to host all events concurrently.
  \item \texttt{assignRooms(events)}: (optional) produce integer room assignments for each event such that no overlapping events share a room.
\end{itemize}

\section{Design overview \& justification}
\subsection{LRU cache — data structures choice}
To achieve \(O(1)\) for both operations:
\begin{itemize}
  \item Use an \texttt{unordered\_map<key, iterator>} for constant-time lookup from key to its node.
  \item Maintain recency with a doubly-linked list (\texttt{std::list<pair<key,value>>}); front = most recent, back = least recent. Given an iterator, \texttt{splice} and removals are constant time.
\end{itemize}
This combination minimizes per-operation cost while keeping implementation straightforward and memory usage proportional to capacity.

\subsection{Scheduler — algorithmic pattern}
\begin{itemize}
  \item \texttt{can\_attend\_all}: sort intervals by start time and check for any start $<$ previous end (strict overlap). Sorting yields \(O(n\log n)\).
  \item \texttt{min\_rooms\_required}: sort by start times, sweep using a min-heap of end times: for each event, free all rooms whose end $\le$ start, then reserve a room (push end). The maximum heap size observed is the required room count — optimal and runs in \(O(n\log n)\).
  \item \texttt{assignRooms}: extend the above by maintaining room ids; reuse freed room ids via a small free-pool (min-heap) so produced ids are compact.
\end{itemize}

\section{Algorithms (high-level pseudocode)}
\subsection{LRU cache}
\begin{tcolorbox}[mybox,title=LRU operations (pseudocode)]
\small
\textbf{get(key)}:
\begin{enumerate}
  \item Lookup node = map[key]. If missing return -1.
  \item Move node to front of list (most recent).
  \item Return node.value.
\end{enumerate}

\textbf{put(key, value)}:
\begin{enumerate}
  \item If key exists: update value, move node to front, return.
  \item If size == capacity: evict node at back (remove from list and map).
  \item Insert new node at front and add map[key] = iterator.
\end{enumerate}
\end{tcolorbox}

\subsection{Scheduler (min rooms)}
\begin{tcolorbox}[mybox,title=Interval partitioning (pseudocode)]
\small
\begin{enumerate}
  \item Sort events by start time.
  \item Initialize empty min-heap of end times.
  \item For each event (s,e):
    \begin{itemize}
      \item While heap not empty and heap.top() $\le$ s: pop() (room freed).
      \item push(e) (occupy a room until e).
      \item track max heap size.
    \end{itemize}
  \item Return tracked max size.
\end{enumerate}
\end{tcolorbox}

\section{Complexity summary}
\begin{table}[h!]
\centering
\begin{tabular}{@{}lcc@{}}
\toprule
Function & Time complexity & Space complexity \\ \midrule
LRU \texttt{get} & \(O(1)\) & \(O(1)\) (per op) \\
LRU \texttt{put} & \(O(1)\) amortized & \(O(1)\) (per op) \\
Scheduler: sort step & \(O(n\log n)\) & \(O(n)\) \\
Scheduler: sweep/heap step & \(O(n\log n)\) total & \(O(n)\) \\ \bottomrule
\end{tabular}
\caption{Asymptotic bounds; \(n\) = number of events.}
\end{table}

\section{Correctness sketch}
\subsection{LRU}
The map ensures exact presence checks and constant-time access; moving nodes to the front on access/update preserves recency ordering. Eviction removes the list's back element, which by construction is the least recently used.

\subsection{Scheduler}
The heap-based greedy is the standard interval partitioning solution: whenever an event starts, any room with end $\le$ start is reusable; if none exist, a new room is required. The algorithm produces the interval depth (maximum concurrency), which equals the minimum number of rooms necessary.

\section{Practical notes and trade-offs}
\begin{itemize}
  \item \textbf{Memory vs speed:} the LRU approach uses per-entry pointers (list nodes) plus hash entries — modest overhead in exchange for \(O(1)\) operations.
  \item \textbf{Alternatives:} language built-ins (Python's \texttt{OrderedDict}) or specialized caches (LRU-K, 2Q) can be used where workload characteristics demand them.
  \item \textbf{Concurrency:} for thread-safety, wrap public operations in a mutex (coarse-grained). For higher throughput, consider lock striping or established concurrent collections.
\end{itemize}

\section{Build, test and run}
\begin{tcolorbox}[mybox]
\textbf{Compile:}
\begin{verbatim}
g++ -std=c++17 main.cpp -O2 -o assignment
\end{verbatim}
\textbf{Run:}
\begin{verbatim}
./assignment
\end{verbatim}
\end{tcolorbox}

\section{Representative test cases}
\begin{itemize}
  \item LRU: insert up to capacity, access older keys to change recency, insert new key to force eviction; assert evicted key is LRU.
  \item Scheduler:
    \begin{itemize}
      \item Non-overlapping adjacent intervals: \((9,10),(10,11)\) $\Rightarrow$ can attend all, rooms = 1.
      \item Overlap: \((9,11),(10,12)\) $\Rightarrow$ cannot attend all, rooms = 2.
      \item Mixed: \((9,10),(9,11),(10,11)\) $\Rightarrow$ rooms = 2.
    \end{itemize}
\end{itemize}

\newpage
\section*{Appendices}
\addcontentsline{toc}{section}{Appendices}

\subsection*{Appendix A — C++: LRU cache}
\addcontentsline{toc}{subsection}{Appendix A — LRU cache}
\lstinputlisting[language=C++,caption={lru\_cache.hpp — unordered\_map + std::list implementation}]{lru_cache.hpp}

\subsection*{Appendix B — C++: Scheduler}
\addcontentsline{toc}{subsection}{Appendix B — Scheduler}
\lstinputlisting[language=C++,caption={scheduler.hpp — canAttendAll, minRoomsRequired, assignRooms}]{scheduler.hpp}

\subsection*{Appendix C — C++: Example main}
\addcontentsline{toc}{subsection}{Appendix C — Example main}
\lstinputlisting[language=C++,caption={main.cpp — minimal test harness}]{main.cpp}

\section*{References}
\begin{enumerate}[leftmargin=*]
  \item E. O'Neil, P. O'Neil, G. Weikum. \textit{The LRU-K Page Replacement Algorithm for Database Disk Buffering}, SIGMOD, 1993.
  \item T. Johnson, D. Shasha. \textit{2Q: A Low Overhead High Performance Buffer Management Replacement Algorithm}, VLDB, 1994.
  \item J. Kleinberg, É. Tardos. \textit{Algorithm Design}, Interval partitioning (lecture notes).
  \item Recent surveys on learned cache replacement (arXiv/conference literature, 2021–2024).
\end{enumerate}



\end{document}